\section{Results}
\label{sec:results}

\subsection{Aggregate Performance}
\label{sec:aggregate}

\Tabref{tab:aggregate} summarizes the main results across all 9 live-salmon papers and 9 dead-salmon papers. \aiscientists demonstrate a clear pattern: reliable detection of strong linear effects, complete failure on interaction effects, and appropriate conservatism on null data.

\begin{table}[t]
    \centering
    \caption{Aggregate results across all 9 live-salmon and 9 dead-salmon papers. \aiscientists reliably detect strong linear predictors but universally miss the interaction term. Best possible rate in \textbf{bold} where achieved.}
    \begin{tabular}{@{}lcc@{}}
        \toprule
        \textbf{Metric} & \textbf{Rate} & \textbf{Papers} \\
        \midrule
        \multicolumn{3}{@{}l}{\textit{Live salmon: true signal detection}} \\
        \quad $X_1$ detected (coef = 2.0) & \textbf{100\%} & 9/9 \\
        \quad $X_3$ detected (coef = $-$1.5) & \textbf{100\%} & 9/9 \\
        \quad $X_2$ detected (coef = 0.5) & 67\% & 6/9 \\
        \quad $X_4 \times X_5$ interaction detected & 0\% & 0/9 \\
        \midrule
        \multicolumn{3}{@{}l}{\textit{Live salmon: error rates}} \\
        \quad Noise features falsely included ($X_6$--$X_{15}$) & \textbf{0\%} & 0/9 \\
        \quad Correlated proxies confused as causal ($X_{16}$--$X_{20}$) & 33\% & 3/9 \\
        \midrule
        \multicolumn{3}{@{}l}{\textit{Dead salmon: null findings}} \\
        \quad Null finding correctly reported & \textbf{100\%} & 9/9 \\
        \quad False positives & \textbf{0\%} & 0/9 \\
        \midrule
        \multicolumn{3}{@{}l}{\textit{Methodological practices (all 18 papers)}} \\
        \quad Multiple comparisons correction mentioned & \textbf{100\%} & 18/18 \\
        \quad Correlation analysis used & \textbf{100\%} & 18/18 \\
        \quad Linear regression used & 89\% & 16/18 \\
        \bottomrule
    \end{tabular}
    \label{tab:aggregate}
\end{table}

\subsection{Linear Effect Detection}
\label{sec:linear}

All 9 live-salmon papers correctly identify $X_1$ and $X_3$ as significant predictors of $Y$.
These features have the strongest marginal correlations with $Y$, consistent with their large true coefficients (2.0 and $-1.5$, respectively).
The weaker predictor $X_2$ (true coefficient 0.5) is detected in only 6 of 9 papers (67\%).
This gap is expected: $X_2$'s correlation with $Y$ is attenuated by the noise floor, making it harder to distinguish from the noise features $X_6$--$X_{15}$.

Detection is consistent across models for the strong effects but varies for $X_2$: each model misses $X_2$ in exactly one of its three runs, suggesting this is near the detection threshold for single-pass correlation-based analysis.

\subsection{Universal Failure on Interaction Effects}
\label{sec:interaction}

The most striking result is that \textbf{zero out of 9 papers detect the $X_4 \times X_5$ interaction term}.
This failure is universal across all three models and all runs.
The interaction contributes $R^2 = 0.083$ when added to a linear-only model---a meaningful effect.
For comparison, $X_2$ has a true coefficient of 0.5 and is detected by 67\% of papers, while the interaction coefficient is 0.8---\textit{larger} in magnitude---yet is never found.

The root cause is methodological: all 9 papers rely on correlation analysis and linear regression, both of which are blind to multiplicative interactions between features.
Detecting $X_4 \times X_5$ would require at least one of: (1) systematic interaction term screening, (2) polynomial feature expansion, (3) residual analysis after linear fitting, (4) tree-based methods, or (5) partial dependence analysis.
No paper attempts any of these approaches.

\subsection{Comparison to Statistical Baselines}
\label{sec:baselines}

\Tabref{tab:baselines} compares the \aiscientist results to standard statistical methods applied to the same data.

\begin{table}[t]
    \centering
    \caption{Statistical baselines on the live salmon dataset. Standard linear methods also miss the interaction, but it is easily detected when explicitly included. $R^2$ values are cross-validated where applicable.}
    \begin{tabular}{@{}lccc@{}}
        \toprule
        \textbf{Method} & \textbf{$R^2$} & \textbf{Features Selected} & \textbf{$X_4 \times X_5$?} \\
        \midrule
        OLS (all features) & 0.77 & All 20 & No \\
        \lassobaseline & --- & $X_1, X_2, X_3, X_4$ & No \\
        OLS + interaction & \textbf{0.85} & $X_1, X_2, X_3, X_4 \times X_5$ & \textbf{Yes} \\
        Random Forest & 0.73 & $X_1, X_3, X_2, X_{16}, X_5, \ldots$ & Implicitly \\
        \midrule
        AI scientists (avg.) & --- & $X_1, X_3$ (100\%); $X_2$ (67\%) & No \\
        \bottomrule
    \end{tabular}
    \label{tab:baselines}
\end{table}

The \aiscientists perform comparably to \lassobaseline and OLS without interaction terms: both approaches identify the linear predictors but miss the nonlinear structure.
The key observation is that the interaction is \textit{easy} to detect when tested---the $R^2$ gap from 0.77 to 0.85 is substantial---but no \aiscientist attempts the diagnostic steps that would reveal it.
A random forest implicitly captures the interaction through recursive partitioning but does not provide an interpretable coefficient.

\subsection{Dead Salmon: Zero False Positives}
\label{sec:dead_salmon}

All 9 dead-salmon papers correctly report that no features have significant relationships with $Y$.
Every paper explicitly states null findings, and all mention or apply multiple comparisons corrections (typically Bonferroni).
This result directly contrasts with \citet{bennett2009neural}, where uncorrected fMRI analysis produced 16 ``significant'' voxels in a dead fish.

The zero false-positive rate suggests that \llms are methodologically conservative: they do not hallucinate statistical significance.
This is a positive finding for the reliability of AI-generated research, though it also reflects the same correlation-based methodology that causes the interaction detection failure---a method that happens to be conservative by construction.

\subsection{Correlated Feature Confusion}
\label{sec:proxy}

In 3 of 9 live-salmon papers (33\%), correlated proxy features ($X_{16}$--$X_{20}$) are listed as significant alongside or instead of the true causal features.
\Tabref{tab:proxy} shows the specific instances.

\begin{table}[t]
    \centering
    \caption{Instances of correlated feature confusion. Three papers include proxy features that are correlated with (but not causal for) $Y$.}
    \begin{tabular}{@{}llp{6cm}@{}}
        \toprule
        \textbf{Model} & \textbf{Run} & \textbf{Proxy Features Included} \\
        \midrule
        \gptmodel & 3 & $X_{17}$ \\
        \claudemodel & 3 & $X_{16}, X_{17}, X_{18}$ \\
        \geminimodel & 2 & $X_{16}, X_{17}, X_{18}$ \\
        \bottomrule
    \end{tabular}
    \label{tab:proxy}
\end{table}

These papers note the correlations between proxy and true features but fail to conclusively distinguish direct from indirect effects.
This is a known challenge even for human analysts: without causal reasoning or regularization that penalizes redundant features, correlated proxies are difficult to separate from true predictors based on observational correlations alone.

\subsection{Coefficient Estimation Accuracy}
\label{sec:coefficients}

When \aiscientists estimate coefficients for the linear terms, their estimates systematically underestimate the true values (\tabref{tab:coef}).

\begin{table}[t]
    \centering
    \caption{Coefficient estimation accuracy. \aiscientists systematically underestimate true coefficients, likely due to reliance on univariate correlations rather than multivariate regression.}
    \begin{tabular}{@{}lccc@{}}
        \toprule
        \textbf{Feature} & \textbf{True Coef.} & \textbf{Mean Estimate} & \textbf{Mean $|$Error$|$} \\
        \midrule
        $X_1$ & 2.0 & 1.72 & 0.28 \\
        $X_2$ & 0.5 & 0.34 & 0.16 \\
        $X_3$ & $-$1.5 & $-$1.38 & 0.12 \\
        \bottomrule
    \end{tabular}
    \label{tab:coef}
\end{table}

The systematic downward bias is consistent with coefficients derived from univariate correlations rather than multivariate regression.
Univariate correlations are attenuated by noise features and multicollinearity from the proxy features ($X_{16}$--$X_{20}$), producing lower apparent effect sizes.
This suggests that even when \aiscientists nominally report running ``regression,'' they may be reasoning from correlation magnitudes rather than fitting a proper multivariate model.

\subsection{Methodological Homogeneity}
\label{sec:methods}

All 9 live-salmon papers employ essentially the same analytical pipeline:
(1) pairwise correlation analysis,
(2) multiple comparisons correction,
(3) multicollinearity assessment, and
(4) linear regression or an approximation thereof.
No paper attempts random forests, gradient boosting, interaction screening, polynomial feature expansion, partial dependence plots, or causal inference methods.

This methodological homogeneity is notable because the three models come from different providers, use different architectures, and were trained on different data.
Despite this diversity, they converge on a single analytical workflow---suggesting that the ``default'' approach to tabular data analysis is deeply embedded in \llm training data and that these models lack the metacognitive capability to recognize when standard approaches are insufficient.
